\chapter{Recommendations and Decision Tree}
\label{chap:final-recommendations}

\section{Mesh and Refinement Recommendations}

For the investigated Molnar geometry with \(l=0.015\,\mathrm{mm}\), the
supervisor-approved mesh roles are H2-PUB for uniform RF--U validation, H1 for
production comparisons, and H0 for development.  The Stage C corrected
locally refined mesh used 10,290 physical elements instead of 12,064 for H1,
a reduction of 14.7\%.  Against a fair four-thread H1 run it reduced wall time
by 16.7\%, CPU time by 14.9\%, and peak memory by 22.2\%.

The refined model reproduced initial stiffness, pre-peak response, and peak
load closely: stiffness differed by 0.061\%, pre-peak NRMSE was 0.089\%, peak
force differed by 0.24\%, and peak displacement was unchanged.  These values
support the workflow for elastic, pre-peak, and peak-response studies within
the tested parameter set.

MISESERI-based pre-refinement successfully identified a useful fracture
corridor, but it did not establish post-peak equivalence.  The post-peak NRMSE
was approximately 24.3\%, and quantitative crack-path comparison detected a
deviation from H1 despite exact repeatability of the refined result.  Coarse
meshes are therefore suitable for workflow development and indicator
generation only when notch geometry, free faces, layer mapping, and the local
\(h/l\) ratio are independently audited.

\section{State-Transfer Recommendations}

For a pre-peak checkpoint transferred to a new nonmatching mesh, use the
bounded D3A3-R4 procedure: complete nodal/integration-point coverage,
source-to-target ordering audits, compatible phase initialization, actual
history transfer, fixed-phase equilibration, KKT evaluation, and a small
release hold with RF and energy discontinuity gates.

The tested one-rank/four-thread D2 continuation matched its serial reference,
but this supports repeatability only for the tested case and is not a general
parallel-scaling result.  ABAQUSER comparison remains unavailable; independent
ODB extraction must be described as the evidence route rather than as an
equivalent ABAQUSER validation.

Further loading cannot assume a fixed active set.  D3D found multiplier
violations at the segment initial state and 3,157 violating nodes by the
endpoint.  A new active-set correction is required when dual admissibility is
lost.  The offline D3D-A1 correction is admissible under frozen F3 history, but
its mechanical restart was not qualified.  It must not be used as an accepted
continuation state.

\section{Execution and Reproducibility Risks}

The closeout exposed three infrastructure risks independent of the scientific
model: platform-dependent text checksums, interpreter-version assumptions, and
tool disappearance after module loading.  Reproducible HPC workflows should
derive hashes in the execution checkout, bind interpreters and compilers by
absolute path, separate login-side repository checks from compute-node
validation, and qualify the environment with a trivial job before a scientific
submission.

\section{Practical Decision Tree}

\begin{description}
  \item[Need only elastic, peak, or pre-peak response?]
  Use the accepted Stage C corrected refined workflow within the tested
  \(l=0.015\,\mathrm{mm}\) and mesh-policy scope.

  \item[Need post-peak crack-path equivalence?]
  Do not use the current Stage C evidence as validation; post-peak RF--U and
  crack-path equivalence are not supported.

  \item[Need transfer to a new mesh at a pre-peak checkpoint?]
  Use the accepted D3A3-R4 bounded transfer and release-hold procedure with
  complete state, ordering, KKT, RF, and energy audits.

  \item[Need continuation after active-set evolution begins?]
  Recompute the active set and independently qualify the corrected mechanical
  restart.  The current D3D-A1 package is a candidate only.

  \item[Need online adaptive remeshing or D3E/post-peak continuation?]
  This is not validated in the thesis and requires a new scientific and
  infrastructure programme.
\end{description}

\section{Final Scope}

The thesis contribution is a bounded, evidence-driven workflow that identifies
where pre-refinement and state transfer are reliable, quantifies their
accuracy/cost benefit, and exposes active-set and execution limitations before
unsupported continuation claims are made.
